% M4_NnaemekaNnachiEgwu.tex
% Priority Inheritance Protocol — Seminar Paper
% Real-Time Systems Seminar, HSHL, Prof. Dr. Stefan Henkler
% Author: Nnaemeka Nnachi-Egwu
% IEEE conference format, target 5 pages

\documentclass[conference]{IEEEtran}

\usepackage{amsmath}
\usepackage{amssymb}
\usepackage{graphicx}
\usepackage{url}
\usepackage{microtype}
\usepackage{booktabs}
\usepackage{array}
\usepackage{multirow}
\usepackage{cite}

\hyphenation{Sched-ul-a-bil-i-ty}

\begin{document}

\title{Priority Inheritance Protocol for Real-Time Resource Synchronization}

\author{\IEEEauthorblockN{Nnaemeka Nnachi-Egwu}
\IEEEauthorblockA{Electronic Engineering\\
Hochschule Hamm-Lippstadt\\
Lippstadt, Germany\\
nnaemeka.nnachi-egwu@stud.hshl.de}}

\maketitle

%---------------------------------------------------------------
\begin{abstract}
Shared resources in real-time systems create the priority inversion
problem: a high-priority task can be blocked by a lower-priority task
for an unbounded duration, invalidating schedulability guarantees.
The Priority Inheritance Protocol (PIP), introduced by Sha, Rajkumar,
and Lehoczky~\cite{Sha1990}, addresses this by dynamically boosting a
resource holder's priority to match that of any blocked task.
This paper presents PIP's formal model, its bounded-inversion
guarantee (Theorem~7.2~\cite{Buttazzo2011}), and its limitations
(deadlock susceptibility, chained blocking). A verified UPPAAL model
and C\texttt{++} simulator demonstrate the protocol's behaviour on a
canonical three-task example. The schedulability analysis uses both
the Liu-Layland utilization test and Response-Time Analysis with
blocking terms. PIP is evaluated against four alternative protocols
and its applicability to embedded systems contexts is assessed.
\end{abstract}

\begin{IEEEkeywords}
priority inversion, priority inheritance, real-time scheduling,
resource access protocols, semaphores, UPPAAL, schedulability analysis
\end{IEEEkeywords}

%---------------------------------------------------------------
\section{Introduction}
\label{sec:intro}

A real-time computer system is one in which the correctness of
behaviour depends not only on the logical results of computation but
also on the \emph{physical instant} at which those results are
produced~\cite{Kopetz2011}.
In this model, a result delivered after its deadline is not merely late;
it is logically incorrect.
Real-time systems must therefore provide timing guarantees under all
operating conditions, including worst-case resource contention.

Shared exclusive resources introduce the \emph{priority inversion}
problem: a high-priority task can be blocked waiting for a resource
held by a lower-priority task, while medium-priority tasks are free
to run, extending the blockage without bound~\cite{Buttazzo2011}.
This violates the guarantee on which rate-monotonic and earliest-deadline-first
scheduling theory is built; the Earliest Deadline First algorithm is
optimal only when there is no resource competition~\cite{Kopetz2011}.

The Priority Inheritance Protocol (PIP)~\cite{Sha1990} bounds this
inversion by letting a resource holder temporarily inherit the priority
of any task it blocks.
This paper explains PIP's design rationale, formal properties, practical
limitations, and its place among the five standard resource access protocols
in Buttazzo Chapter~7~\cite{Buttazzo2011}.

%---------------------------------------------------------------
\section{Background: Task Model and Problem}
\label{sec:background}

\subsection{Task Parameters}

Following~\cite{Buttazzo2011} §7.3, each periodic task $\tau_i$ is
characterised by: nominal priority $P_i$ (assigned by Rate-Monotonic
or Deadline-Monotonic rules, lower number = higher priority),
worst-case computation time $C_i$, period $T_i$, and relative
deadline $D_i$.
Active (dynamic) priority $p_i$ may deviate from $P_i$ under a
resource access protocol.

\subsection{Semaphores and Critical Sections}

Exclusive resources are protected by binary semaphores.
A critical section $z_{i,k}$ of task $\tau_i$ on semaphore $S_k$ has
duration $\delta_{i,k}$; the collection of all its critical sections
is $Z_i = \{z_{i,k}\}$.
The \texttt{wait}($S_k$) / \texttt{signal}($S_k$) primitives enforce mutual
exclusion.
A task blocked on a busy semaphore moves to the WAITING state;
the three-state model (READY, RUNNING, WAITING) is taught in the
course slides (RTS Introduction, slide~31).

\subsection{Priority Inversion}

Without a resource access protocol, consider three tasks
$\tau_1 \succ \tau_2 \succ \tau_3$ in priority (Fig.~7.4
in~\cite{Buttazzo2011}):

\begin{enumerate}
  \item $\tau_3$ acquires semaphore $S_R$ and enters its critical section.
  \item $\tau_1$ (highest priority) is released and immediately preempts
        $\tau_3$. $\tau_1$ calls \texttt{wait}($S_R$) but $S_R$ is held
        by $\tau_3$, so $\tau_1$ blocks.
  \item $\tau_2$ (which does not use $S_R$) preempts $\tau_3$.
  \item $\tau_1$ remains blocked for the entire duration of $\tau_2$'s
        execution, even though $\tau_2$ has lower priority than $\tau_1$.
\end{enumerate}

If any number of medium-priority tasks arrive during this window, the
blocking of $\tau_1$ is \emph{unbounded}. This is
the priority inversion problem.

%---------------------------------------------------------------
\section{Protocol Landscape}
\label{sec:landscape}

Table~\ref{tab:protocols} summarises the five resource access protocols
in Buttazzo Chapter~7~\cite{Buttazzo2011}.

\begin{table*}[t]
\caption{Resource access protocol comparison (Buttazzo Table 7.5, p.~248)}
\label{tab:protocols}
\centering
\small
\begin{tabular}{@{}lccccc@{}}
\toprule
Property         & NPP & HLP & PIP & PCP & SRP \\
\midrule
Deadlock-free    & \checkmark & \checkmark & \textbf{--} & \checkmark & \checkmark \\
Chained blocking & -- & -- & \textbf{Yes} & -- & -- \\
Transparent      & \checkmark & -- & \checkmark & -- & -- \\
Max blockings    & 1 & 1 & $\alpha_i$ & 1 & 1 \\
Implementation   & Easy & Easy & \textbf{Hard} & Medium & Easy \\
Dynamic prio.    & -- & -- & -- & -- & \checkmark \\
\bottomrule
\end{tabular}
\end{table*}

\textbf{NPP} (Non-Preemptive Protocol): tasks run non-preemptively
in critical sections; single blocking, but over-conservative.
\textbf{HLP/IPC} (Highest Locker Priority / Immediate Priority
Ceiling): tasks inherit the ceiling of the semaphore at entry time;
single blocking but requires ceiling declarations.
\textbf{PIP}: defers blocking to actual access time, no code change
required (transparent), bounded inversion, but no deadlock prevention.
\textbf{PCP} (Priority Ceiling Protocol, also~\cite{Sha1990}):
adds lock-grant rule eliminating deadlock and chained blocking.
\textbf{SRP} (Stack Resource Policy~\cite{Baker1991}): extends PCP
to multi-unit resources and EDF scheduling, enables stack sharing.

%---------------------------------------------------------------
\section{Priority Inheritance Protocol}
\label{sec:pip}

\subsection{Algorithm}

PIP~\cite{Sha1990} is defined via two primitives (Buttazzo §7.6.4,
pp.~224--225), reproduced here with slight notation adjustments:

\medskip
\noindent\textbf{pi\_wait}($S_k$):
\begin{enumerate}
  \item If $S_k$ is free, lock it and assign $S_k.\mathit{holder} \leftarrow \tau_{\mathit{exe}}$.
  \item Otherwise, add $\tau_{\mathit{exe}}$ to $S_k$'s wait queue and
        set $\tau_{\mathit{exe}}$ to BLOCKED.
        Let $\tau_j = S_k.\mathit{holder}$; apply inheritance:
        \begin{equation}
          p_j(R_k) = \max\!\bigl\{P_j,\;\max_h\{P_h \mid \tau_h\text{ blocked on }R_k\}\bigr\}.
          \label{eq:pip}
        \end{equation}
        If $\tau_j$ is itself blocked on another semaphore $S_m$,
        apply the rule transitively to $S_m.\mathit{holder}$.
\end{enumerate}

\noindent\textbf{pi\_signal}($S_k$):
\begin{enumerate}
  \item Unlock $S_k$.
  \item Wake the highest-priority task in $S_k$'s wait queue (if any)
        and grant it the lock.
  \item Restore $\tau_j$'s active priority to
        $\max\bigl\{P_j,\,\max\{p_h \mid \tau_h \text{ still blocked by } \tau_j\}\bigr\}$.
  \item If $p_j < p^*$ (highest ready task), perform a context switch.
\end{enumerate}

Inheritance is \emph{transitive}: if $\tau_3$ blocks $\tau_2$ which
blocks $\tau_1$, then $\tau_3$ inherits $P_1$ via $\tau_2$
(§7.6.1, p.~215, Lemma~7.2 requires nested critical sections for
transitivity to arise).

\subsection{Push-Through Blocking}

A key consequence of PIP is \emph{push-through blocking}
(Buttazzo Lemma~7.1, p.~218): if $\tau_j$ (low priority) holds
$S_R$ and blocks $\tau_i$ (high priority), then $\tau_j$ inherits
$P_i$ and thereby prevents any intermediate-priority task from preempting
it.
Lemma~7.1 identifies the exact condition:
semaphore $S_k$ causes push-through blocking to $\tau_i$ if and only
if there exists a task with priority greater than $P_i$ using $S_k$
and a task with priority less than $P_i$ using $S_k$.

\subsection{Blocking Time Bound}

Let $l_i$ be the number of distinct lower-priority tasks that can block
$\tau_i$, and $s_i$ the number of distinct semaphores accessed by both
$\tau_i$ and lower-priority tasks.
Then (Theorem~7.2~\cite{Sha1990}, Buttazzo p.~219):
\begin{equation}
  \alpha_i = \min(l_i, s_i).
  \label{eq:alpha}
\end{equation}
$\tau_i$ is blocked by at most $\alpha_i$ critical sections.
The worst-case blocking time is computed using the algorithm in
Fig.~7.11~\cite{Buttazzo2011} ($\mathcal{O}(mn^2)$):
\begin{align}
  B_i^l &= \sum_{j>i}\sum_k \delta_{j,k} \cdot \mathbb{1}[\text{block}]
  \label{eq:bl}\\
  B_i^s &= \text{(sum by semaphore)}
  \label{eq:bs}\\
  B_i   &= \min(B_i^l, B_i^s).
  \label{eq:bi}
\end{align}
This bound is \emph{not tight}: $B_i^l$ may sum critical sections
guarded by the same semaphore, which Lemma~7.4 shows cannot both block
simultaneously (p.~222).

\subsection{Schedulability}

With blocking terms, the Liu-Layland sufficient test (Buttazzo Eq.~7.19)
becomes:
\begin{equation}
  \forall i:\quad \sum_{h<i}\frac{C_h}{T_h} + \frac{C_i + B_i}{T_i}
  \;\leq\; i\!\left(2^{1/i}-1\right).
  \label{eq:ll}
\end{equation}
When this test is inconclusive, the exact Response-Time Analysis
(Eq.~7.22) is applied iteratively:
\begin{equation}
  R_i = C_i + B_i + \sum_{h<i}\!\left\lceil\frac{R_i}{T_h}\right\rceil C_h.
  \label{eq:rta}
\end{equation}

\subsection{Limitations}

\textbf{No deadlock prevention.} If two tasks acquire two semaphores
in reverse order, a lock cycle can form (Fig.~7.13~\cite{Buttazzo2011},
p.~226): PCP's lock-grant rule is needed to prevent this.

\textbf{Chained blocking.} If $\tau_i$ accesses $n$ distinct
semaphores each held by a different lower-priority task, it can be blocked
$n$ times (Fig.~7.12, p.~225); the number is bounded by $\alpha_i$
from (\ref{eq:alpha}), but PCP limits this to exactly one.

\textbf{Kernel complexity.} The transitive chain traversal in
\texttt{pi\_wait} adds latency inside an atomic kernel section;
Buttazzo rates PIP's implementation as ``hard'' --- the only protocol
with that rating (Table~7.5, p.~248).

%---------------------------------------------------------------
\section{UPPAAL Verification}
\label{sec:uppaal}

\subsection{Model Structure}

A UPPAAL timed-automata model (based on~\cite{Behrmann2004}) was
constructed with four components:
(1)~a parametric \texttt{Task} template instantiated for $\tau_1$,
$\tau_2$, $\tau_3$, each with locations \textsc{Idle}, \textsc{Ready},
\textsc{Running}, \textsc{CriticalSection}, \textsc{Blocked};
(2)~a \texttt{ResourceManager} automaton with locations
\textsc{Free}, \textsc{Held}, and committed \textsc{Granting},
encoding the binary semaphore $S_R$ with PIP inheritance logic
(Eq.~\ref{eq:pip});
(3)~an \texttt{Observer} process for property monitoring.

Shared global variables encode active priorities
(\texttt{int[1,9]~act\_prio[3]}), the semaphore holder
(\texttt{int[-1,2]~holder}), blocking duration counters
(\texttt{int[0,20]~block\_dur[3]}), and observer flags
(\texttt{bool~blocking\_tau1}).
Synchronisation channels \texttt{lock\_req[N]!}, \texttt{lock\_ack[N]!},
\texttt{unlock\_rel[N]!} coordinate lock operations between tasks and
the resource manager.
The key inheritance transition in \texttt{ResourceManager} fires on
\texttt{lock\_req[0]?} while \texttt{holder==2} (i.e., $\tau_3$ holds
the lock when $\tau_1$ requests it), and immediately applies the
assignment
\begin{quote}\small\ttfamily
act\_prio[2] := (P1 < act\_prio[2]) ? P1 : act\_prio[2]
\end{quote}
implementing Eq.~(\ref{eq:pip}) atomically.

\subsection{Properties Verified}

The five verification queries target distinct PIP correctness properties.

\textbf{P1 (Deadlock-free):}
\begin{quote}\small\ttfamily
A[] not deadlock
\end{quote}
No lock cycle can form with a single resource and consistent acquisition order;
this query would \emph{fail} if two resources were acquired in reverse order
(Fig.~7.13~\cite{Buttazzo2011}).

\textbf{P2 (No unbounded inversion):}
\begin{quote}\small\ttfamily
A[] (blocking\_tau1 imply not\\
\quad(tau2.Running and not blocked\_tau2))
\end{quote}
When $\tau_1$ is blocked on $S_R$, $\tau_2$ cannot run freely — verifying
push-through blocking (Lemma~7.1). An \emph{invariant}: must hold in every
reachable state.

\textbf{P3 (Blocking within bound):}
\begin{quote}\small\ttfamily
A[] (block\_dur[0]<=4 and block\_dur[1]<=4\\
\quad and block\_dur[2]==0)
\end{quote}
Observed blocking stays within $B_1{=}4$, $B_2{=}4$, $B_3{=}0$, linking
the formal model to the schedulability analysis.

\textbf{P4 (Inheritance reachable):}
\begin{quote}\small\ttfamily
E<> (act\_prio[2]==1)
\end{quote}
There exists a reachable state where $\tau_3$ holds inherited priority~$P_1$,
confirming the inheritance mechanism fires (non-vacuous witness).

\textbf{P5 (Priority restored):}
\begin{quote}\small\ttfamily
A[] (rm.Free imply act\_prio[2]==3)
\end{quote}
Whenever $S_R$ is free, $\tau_3$'s priority equals~$P_3$, verifying
\texttt{restore\_prio} correctly undoes every inheritance.

%---------------------------------------------------------------
\section{Application Example}
\label{sec:example}

\subsection{Task Set}

Three periodic tasks are scheduled by Rate-Monotonic (shorter period =
higher priority):
$\tau_1$~($C{=}1, T{=}6, D{=}6$, prio~1, uses~$S_R$, CS~len~1),
$\tau_2$~($C{=}2, T{=}8, D{=}8$, prio~2, no shared resource),
$\tau_3$~($C{=}4, T{=}24, D{=}24$, prio~3, uses~$S_R$, CS~len~4).
The semaphore ceiling $C(S_R) = P_1$ (highest priority using $S_R$).

\subsection{Priority Inversion Trace (Without PIP)}

$\tau_3$ is released at $t{=}0$ and acquires $S_R$ immediately;
at $t{=}2$, $\tau_1$ and $\tau_2$ are released.
$\tau_1$ preempts $\tau_3$, calls \texttt{pi\_wait}($S_R$), and blocks
($S_R$ held by $\tau_3$, which retains priority~$P_3$).
$\tau_2$ (priority~$P_2$) preempts $\tau_3$ at $t{=}2$.
$\tau_1$ remains blocked for the entire duration of $\tau_2$'s
execution (2~units), plus $\tau_3$'s remaining CS (2~units):
\emph{total blocking = 4~units}, of which 2 units are due to the
lower-priority task~$\tau_2$ --- priority inversion.
Table~\ref{tab:trace} shows the schedule.

\begin{table*}[t]
\caption{Schedule trace (one column per time unit). W = waiting for $S_R$;
         CS = in critical section; $\dagger$ = inherits $P_1$.}
\label{tab:trace}
\centering
\small
\setlength{\tabcolsep}{3pt}
\begin{tabular}{@{}l|ccccccccc@{}}
\toprule
Task & 0 & 1 & 2 & 3 & 4 & 5 & 6 & 7 & \ldots \\
\midrule
\multicolumn{10}{@{}l}{\emph{Without PIP:}} \\
$\tau_1$ & -- & -- & W & W & W & W & \textbf{run} & -- & \\
$\tau_2$ & -- & -- & run & run & -- & -- & -- & -- & \\
$\tau_3$ & CS & CS & -- & -- & CS & CS & -- & -- & \\
\midrule
\multicolumn{10}{@{}l}{\emph{With PIP:}} \\
$\tau_1$ & -- & -- & W & W & \textbf{run} & -- & -- & -- & \\
$\tau_2$ & -- & -- & rdy & rdy & run & run & -- & -- & \\
$\tau_3$ & CS & CS & CS$^\dagger$ & CS$^\dagger$ & -- & -- & -- & -- & \\
\bottomrule
\end{tabular}
\end{table*}

\subsection{PIP Trace (With PIP)}

At $t{=}2$, when $\tau_1$ blocks on $S_R$, PIP applies
Eq.~(\ref{eq:pip}): $p_3(S_R) = \max(P_3, P_1) = P_1 = 1$.
$\tau_3$ now has priority~1; $\tau_2$ (priority~2) \emph{cannot}
preempt it.
$\tau_3$ runs uninterrupted until it signals $S_R$ at $t{=}4$
($2$ remaining CS units).
$\tau_3$'s priority is restored to $P_3$.
$\tau_1$ is unblocked and runs to completion.
\emph{Blocking = 2~units} --- bounded by $\tau_3$'s remaining CS,
with no contribution from $\tau_2$.

\subsection{Blocking Time Computation}

With $\delta_{3,R} = 4$ (CS length, continuous-time convention;
Buttazzo's discrete formula subtracts~1, giving $\delta_{3,R}{-}1=3$):

\noindent\textbf{$B_1$ (direct blocking):}
$\gamma_1 = \{Z_{3,R}\}$ since $P_3 < P_1$ and $C(S_R){=}P_1 \geq P_1$
(Lemma~7.5).
$l_1 = s_1 = 1$, so $\alpha_1 = 1$.
$B_1^l = B_1^s = \delta_{3,R} = 4$, hence $B_1 = 4$.

\noindent\textbf{$B_2$ (push-through blocking, Lemma~7.1):}
$S_R$ is accessed by $\tau_1$ ($P_1 > P_2$) and $\tau_3$ ($P_3 < P_2$),
so $\tau_3$ can inherit $P_1 > P_2$, blocking $\tau_2$.
$l_2 = s_2 = 1$, $\alpha_2 = 1$, $B_2 = 4$.

\noindent\textbf{$B_3 = 0$} ($\tau_3$ is lowest-priority).

\begin{align*}
  B_1 &= \min(B_1^l, B_1^s) = \min(4, 4) = 4 \\
  B_2 &= \min(4, 4) = 4 \\
  B_3 &= 0
\end{align*}

\subsection{Schedulability Analysis}

$U = \tfrac{1}{6} + \tfrac{2}{8} + \tfrac{4}{24} = 0.583$.

Liu-Layland (Eq.~\ref{eq:ll}) is inconclusive for $\tau_2$
($0.917 > 0.828$; sufficient test only).
Response-Time Analysis (Eq.~\ref{eq:rta}) gives exact results:
\begin{align*}
  R_1 = 5 &\leq D_1 = 6,\\
  R_2 = 8 &\leq D_2 = 8,\\
  R_3 = 8 &\leq D_3 = 24.
\end{align*}
All tasks are schedulable under RM with PIP.
This illustrates Buttazzo's observation that the Liu-Layland test
becomes only sufficient in the presence of blocking terms (p.~246).

%---------------------------------------------------------------
\section{C\texttt{++} Simulator}
\label{sec:impl}

A discrete-time simulator was implemented in C\texttt{++}17 to
validate the hand trace and confirm that blocking bounds hold
over a full hyperperiod ($\mathrm{LCM}(6,8,24){=}24$ time units),
matching the same task parameters as Section~\ref{sec:example}.

\subsection{Design}

The core data structures mirror Buttazzo's notation directly.
Each \texttt{Task} struct holds \texttt{nom\_prio} (nominal, fixed),
\texttt{act\_prio} (active, modified by PIP), \texttt{remaining}
(computation time left), CS state variables (\texttt{has\_lock},
\texttt{cs\_rem}), and blocking metadata (\texttt{blk\_start},
\texttt{max\_blk}).
The single \texttt{Semaphore} struct tracks \texttt{holder} and
a wait queue of task indices.

The \texttt{inherit\_up(holder, req\_prio)} function implements
Eq.~(\ref{eq:pip}) recursively: if the requested priority is higher
than the holder's current active priority, boost the holder and check
whether the holder is itself in a wait queue, triggering another
\texttt{inherit\_up} call on that queue's holder (transitive
inheritance, Lemma~7.2).
\texttt{restore\_prio} scans all semaphore wait queues on release
and sets the task's active priority to the maximum of its nominal
priority and the highest priority among remaining waiters.

\subsection{Simulation Results}

\textbf{Without PIP:} At $t{=}4$, $\tau_1$'s second instance preempts
$\tau_3$ and blocks on $S_R$.
Since $\tau_3$ retains its nominal priority~$P_3$, $\tau_2$ preempts
$\tau_3$ when it becomes ready.
$\tau_1$ remains blocked for 5~units (inversion window:
$\tau_2$ executes, then $\tau_3$ finishes CS).
This exceeds the bound $B_1{=}4$ --- consistent with the theory:
without a protocol there is no bound.
$\tau_1$ misses its deadline ($R_1{=}7 > D_1{=}6$).

\textbf{With PIP:} At $t{=}4$, $\tau_1$ blocks and \texttt{inherit\_up}
immediately raises $\tau_3$'s active priority to $P_1{=}1$.
$\tau_2$ (priority~$P_2{=}2$) cannot preempt $\tau_3$-at-$P_1$.
$\tau_3$ completes its CS at $t{=}7$; \texttt{restore\_prio} sets
$\tau_3$ back to~$P_3{=}3$.
$\tau_1$'s observed blocking: $7{-}4{=}3 \leq B_1{=}4$~\checkmark.
All deadlines met across the full hyperperiod.
Simulator output: \texttt{tau1: observed=3, B1\_bound=4~[OK]}.

%---------------------------------------------------------------
\section{Critical Evaluation}
\label{sec:eval}

\subsection{Strengths}

\textbf{Bounded inversion.} Theorem~7.2 provides a formal upper bound
$\alpha_i = \min(l_i, s_i)$ on the number of critical sections
blocking $\tau_i$; unbounded priority inversion is impossible under PIP.

\textbf{Transparency.} PIP requires no change to task code
(§7.10, p.~247); \texttt{pi\_wait}/\texttt{pi\_signal} are drop-in
replacements for standard semaphore primitives.
This makes PIP the natural choice for POSIX systems
(\texttt{PTHREAD\_PRIO\_INHERIT}~\cite{Sha1990}) and legacy RTOS migration.
The 1997 Mars Pathfinder anomaly was resolved by enabling
\texttt{PTHREAD\_PRIO\_INHERIT} on a VxWorks mutex --- a one-line
configuration change, no code restructuring.

\textbf{Schedulability integration.} The $B_i$ values feed directly
into standard schedulability tests (Eqs.~\ref{eq:ll}--\ref{eq:rta}).

\subsection{Limitations}

\textbf{No deadlock prevention.} Two tasks acquiring two semaphores
in reverse order create a lock cycle that PIP cannot detect or break
(Fig.~7.13~\cite{Buttazzo2011}). The fix is either semaphore
total-ordering (by ID) or replacement with PCP.

\textbf{Chained blocking.} Multiple sequential semaphore accesses
can produce $\alpha_i > 1$ blockings (Fig.~7.12). PCP eliminates
this by guaranteeing a single blocking per task.

\textbf{WCET input sensitivity.} Equations~\ref{eq:bl}--\ref{eq:bi}
require exact $\delta_{j,k}$ values (CS WCETs).
Abella et al.~\cite{Abella2015} show that none of the current WCET
analysis methods is fully trustworthy on modern hardware;
under-estimated WCETs may make the system appear schedulable when it
is not.

\textbf{Uniprocessor scope.} Buttazzo Chapter~7 explicitly addresses
uniprocessor systems (§7.3).
On multiprocessors, priority boosting requires inter-processor
interrupts (IPIs) with non-deterministic latency.
Gracioli et al.~\cite{Gracioli2013} measure such overhead in a real
RTOS and show it has no equivalent in Buttazzo's blocking formula;
PIP as defined by~\cite{Sha1990} does not directly transfer to
multiprocessor contexts.

\subsection{Transfer Analysis}

PIP is suitable for: single-core RTOS with few shared resources and
total semaphore ordering enforced; POSIX-based systems requiring
transparent legacy migration; automotive ECUs on classic AUTOSAR
single-core platforms.
PIP is unsuitable for: multiprocessor/multicore systems; EDF-scheduled
systems (use SRP~\cite{Baker1991}); safety-critical contexts requiring
certified deadlock freedom (use PCP).

%---------------------------------------------------------------
\section{Conclusion}
\label{sec:conclusion}

The Priority Inheritance Protocol solves the priority inversion problem
by dynamically elevating resource holders' priorities, bounding
$\tau_i$'s blocking to at most $\alpha_i = \min(l_i, s_i)$ critical
sections (Theorem~7.2).
Its transparency allows deployment without code changes, making it
attractive for legacy systems and POSIX environments.
However, PIP does not prevent deadlocks or chained blocking, both
eliminated by the Priority Ceiling Protocol at the cost of
transparency.
The exact Response-Time Analysis (Eq.~\ref{eq:rta}) is necessary for
schedulability decisions when the Liu-Layland test is inconclusive.
On multiprocessors, PIP's assumptions break due to IPI latency; this
remains an active research direction.

\section*{Acknowledgements}

I thank Prof.\ Dr.\ Stefan Henkler for introducing the topic of
real-time resource synchronisation and for the structured seminar
milestones that guided this work.
AI usage during this work is documented in the accompanying AI usage
protocol files as required by the seminar protocol.

\bibliographystyle{IEEEtran}
\bibliography{references}

\end{document}
